% Part: first-order-logic
% Chapter: lindstrom
% Section: ls-property

\documentclass[../../../include/open-logic-section]{subfiles}

\begin{document}

\olfileid{mod}{lin}{lsp}

\olsection{Compactness and L\"owenheim--Skolem Properties}

我们现在将紧致性与 L\"owenheim--Skolem 性质自然地推广到抽象逻辑的情形。

\begin{defn}
若对于抽象逻辑 $\tuple{L, \models_L}$，只要 $L(\Lang{L})$-语句集合 $\Gamma$ 的每个有限子集 $\Gamma_0 \subseteq \Gamma$ 可满足，$\Gamma$ 即可满足，则称其具有\emph{紧致性}。
\end{defn}

\begin{defn}
若 $\tuple{L, \models_L}$ 的任意可满足的 $\Gamma$ 都有一个可数模型，则称其具有\emph{向下 L\"owenheim--Skolem 性质}。
\end{defn}

定义于 \olref[bas][pis]{defn:partialisom} 的“部分同构”概念是纯“代数”的（即，其定义不依赖于语言的语句，而只依赖于结构所属语言~$\Lang{L}$ 提供的常项），因此它适用于抽象逻辑的情形。对于一阶逻辑，由 \olref[bas][pis]{thm:p-isom2} 可知，若两个结构是部分同构的，则它们初等等价。这一证明不能直接推广到抽象逻辑，因为对任意 $!E \in L(\Lang{L})$，公式归纳法未必可用；然而，只要 L\"owenheim--Skolem 性质成立，该定理依然成立。

\begin{thm}
\ollabel{thm:abstract-p-isom}
设 $\tuple{L, \models_L}$ 是一个具有 L\"owenheim--Skolem 性质的正规逻辑。则任何两个部分同构的结构在 $\tuple{L, \models_L}$ 意义下都是初等等价的。
\end{thm}

\begin{proof}
假设 $\Struct{M} \simeq_p \Struct{N}$，但对某个 $!E$ 有 $\Struct{M} \models_L !E$ 而 $\Struct{N} \not\models_L !E$。根据同构性质，我们可以假定 $\Domain{M}$ 与 $\Domain{N}$ 不相交；根据扩张性质，我们可以假定 $!E$ 属于某个有限语言 $\Lang{L}$ 的 $L(\Lang{L})$。令 $\mathcal{I}$ 为 $\Struct{M}$ 与 $\Struct{N}$ 之间的部分同构集，且不失一般性，假定：若 $p \in \mathcal{I}$ 且 $q \subseteq p$，则 $q \in \mathcal{I}$。

$\Domain{M}^{<\omega}$ 是 $\Domain{M}$ 中元素的有穷序列的集合。令 $S$ 为 $\Domain{M}^{<\omega}$ 上表示连接的三元关系，即对于 $\mathbf{a}, \mathbf{b},
\mathbf{c} \in \Domain{M}^{<\omega}$，$S(\mathbf{a}, \mathbf{b},
\mathbf{c})$ 成立当且仅当 $\mathbf{c}$ 是 $\mathbf{a}$ 和 $\mathbf{b}$ 的连接；并令 $T$ 为三元关系，使得对于 $b \in M$ 与 $\mathbf{a}, \mathbf{c} \in \Domain{M}^{<\omega}$，$T(\mathbf{a}, b, \mathbf{c})$ 成立当且仅当 $\mathbf{a} =
a_1, \dots a_n$ 且 $\mathbf{c} = a_1, \dots a_n, b$。取新的三元谓词符号 $P$ 与 $Q$，构成结构 $\Struct{M}^*$，其论域为 $\Domain{M} \cup \Domain{M}^{<\omega}$，以 $\Struct{M}$ 为子结构，并用连接关系 $S$ 与 $T$ 解释 $P$ 与 $Q$（因此 $\Struct{M}^*$ 的语言为 $\Lang{L} \cup \{ P, Q\}$）。

类似地定义 $\Domain{N}^{<\omega}$、$S'$、$T'$、$P'$、$Q'$ 和 $\Struct{N}^*$。根据假设$\Struct{M} \simeq_p
\Struct{N}$，在 $\Domain{M}^{<\omega}$ 与 $\Domain{N}^{<\omega}$ 之间存在一个关系 $I$，使得 $I(\mathbf{a}, \mathbf{b})$ 成立当且仅当 $\mathbf{a}$ 和 $\mathbf{b}$ 是同构的且满足 \olref[bas][pis]{defn:partialisom} 中的来回条件。现在，令 $\Struct{M}$ 为这样一个结构：其论域是 $\Struct{M}^*$ 与 $\Struct{N}^*$ 的论域之并，以 $\Struct{M}^*$ 与 $\Struct{N}^*$ 为子结构，其语言包含一个由关系 $I$ 解释的额外二元谓词符号 $R$，以及表示论域 $\Domain{M}^*$ 与 $\Domain{N}*$ 的谓词符号。

\begin{figure}[h]
  \centering
  \begin{tikzpicture}[node distance=2cm, auto, thick, >=stealth']
    \draw [rounded corners] (0,0) -- (8,0) -- (8,4) -- (0,4) --  cycle;
    \draw (2,2) circle (0.5cm);
    \draw (2,2) circle (1.25cm);
    \draw (6,2) circle (0.5cm);
    \draw (6,2) circle (1.25cm);
    \path node at (0.75,3.5) {\large $\Struct{M}$};
    \path node at (2,2) {\large $\Struct{M}$};
    \path node at (6,2) {\large $\Struct{N}$};
    \path node at (3.5,1) {\large $\Struct{M}^*$};
    \path node at (7.5,1) {\large $\Struct{N}^*$};
    \node (Idom) at (2.8,2) {};
    \node (Irng) at (5.2,2) {};
    \draw[<->, bend left] (Idom) to node {\large $I$} (Irng) ;
  \end{tikzpicture}
  \caption{结构 $\Struct{M}$ 及其内部的部分同构。}
\end{figure}

关键的观察是：在结构 $\Struct{M}$ 的语言中，存在一个在 $\Struct{M}$ 中为真的\emph{一阶}语句 $!D_1$，断言 $\Struct{M} \models_L !E$ 且 $\Struct{N} \not\models_L !E$（这需要相对化性质）；也存在一个在 $\Struct{M}$ 中为真的\emph{一阶}语句 $!D_2$，断言 $\Struct{M} \simeq_p \Struct{N}$ 经由部分同构 $I$ 成立。根据 Löwenheim--Skolem 性质，$!D_1$ 与 $!D_2$ 在一个可数模型 $\Struct{M}_0$ 中同时为真，该模型包含部分同构的子结构 $\Struct{M}_0$ 与 $\Struct{N}_0$，使得 $\Struct{M}_0 \models_L !E$ 与 $\Struct{N}_0
\not\models_L !E$。但是，根据 \olref[bas][pis]{thm:p-isom1}，可数的部分同构结构实际上是同构的，这与正规抽象逻辑的同构性质相矛盾。
\end{proof}

\end{document}